
\subsubsection{6.1 Support for the Causal Chain}

The primary causal chain (feature complexity → gradient asymmetry → temporal gap) is supported with varying degrees of evidence. The complexity hierarchy (H-M2) is established on all three metrics with uncorrected significance, though only one metric survives Bonferroni correction. The gradient asymmetry (H-M1) is quantitatively large (GDR=6.977, approximately 7× ratio) and consistent across seeds, but formal statistical significance is not achievable with the planned Wilcoxon test at n=3. The temporal gap (H-E1) is the most formally supported result: p=0.022, t-statistic=4.619 across three seeds, with consistent gap area.

These results are consistent with the Frequency Principle \citep{Xuetal2019} and Simplicity Bias \citep{Shahetal2020}, providing what may be the first quantitative operationalization of these theoretical predictions on a standard spurious correlation benchmark. However, the proof-of-concept scope (30 epochs, one dataset, one architecture) limits the strength of claims that can be drawn.

\subsubsection{6.2 Persistence of Gradient Dominance}

GDR > 1.0 persists throughout the entire 30-epoch training run, not only in early epochs. This is inconsistent with the view that spurious gradient dominance is transient and resolves at t\\textit{. The closing of δ(t) at t\} reflects increasing core probe accuracy (core features being encoded) rather than a decrease in spurious gradient magnitude. This has implications for gradient-based intervention design: suppressing spurious gradients at t\* alone may be insufficient if GDR remains elevated.

\subsubsection{6.3 DFR Robustness and the Role of ImageNet Pretraining}

The most notable unexpected finding is that DFR achieves WGA=0.806 at epoch 1, before any Waterbirds-specific training has substantially modified the ImageNet-pretrained representation. The original hypothesis predicted that DFR efficacy would increase with post-t\* training depth. The data are inconsistent with this prediction.

Three potential explanations are noted:

\begin{enumerate}
\item ImageNet pretraining provides a feature floor sufficient for DFR's class-balanced reweighting to recover core signal, making the degree of Waterbirds-specific training of secondary importance.
\item DFR's class-balanced logistic regression may suppress spurious correlations inherently by averaging over a balanced class distribution, independent of backbone quality.
\item Bird-texture feature disentanglement in the pretrained ResNet-50 representation may be sufficient for DFR even with minimal fine-tuning.
\end{enumerate}

Distinguishing these explanations would require controlled ablations (e.g., scratch-trained backbones, randomized ImageNet weights). The observation that DFR WGA improves only modestly from epoch 1 (0.806) to epoch 30 (0.871) while ERM WGA rises from 0.217 to 0.730 is consistent with explanation 1, but correlational evidence cannot confirm causation.

\subsubsection{6.4 Seed 3 and t\* = 0}

Seed 3 exhibits t\\textit{=0, meaning the primary threshold (δ < 0.02 for 3 consecutive checkpoints) was met at the first checkpoint. This does not indicate absence of a spurious-before-core dynamic; the gap area for seed 3 is positive (A=0.021). The t\}=0 outcome reflects initialization-dependent convergence behavior and motivates reporting CI alongside mean and std rather than relying on a single-seed estimate.

\subsubsection{6.5 Limitations}

\textbf{L1 — 30-epoch proof-of-concept.} All results are from a 30-epoch training run. Under standard 300-epoch training, t\\textit{ would likely occur proportionally later, yielding a longer gap window and a wider condition spread for H-M4. Quantitative claims about window duration and t\} timing should be treated as PoC estimates.

\textbf{L2 — CelebA replication not achieved.} Network access restrictions blocked CelebA download. Cross-dataset generalizability of the δ(t) framework is not confirmed.

\textbf{L3 — H-M1 Wilcoxon underpowered.} With n=3 early checkpoints, scipy.stats.wilcoxon achieves a minimum p-value of 0.125. Formal gradient asymmetry significance requires an extended window (n≥6 checkpoints) or a different test.

\textbf{L4 — H-M4 metric confound.} DFR improvement = DFR WGA − ERM WGA is confounded by the ERM-WGA ceiling effect. Redesigning H-M4 with DFR absolute WGA as the dependent variable, or using partial correlation controlling for ERM WGA, is deferred to future work.

\textbf{L5 — Quadrant patch extraction.} Segmentation masks were not available in the dataset copy used for H-M2. Quadrant-based extraction (top 40\% background, center 60\% foreground) was used throughout, introducing patch impurity. Despite this, all three complexity metrics show the predicted direction.

\textbf{L6 — Single architecture and pretraining.} All results use ResNet-50 with ImageNet pretraining. Results for scratch-trained models or transformer architectures may differ, particularly the epoch-1 DFR WGA finding.

